\section{Geometric data and admissible preparation}\label{sec:preparation}

\subsection{Complete cylinder and maximal domains}
Let $\K=L^2(S^2,\Sigma;\dd\Omega)$ and
$\HH=L^2(\R,\dd v;\K)$. Let $D$ be the self-adjoint unit-sphere Dirac
operator, $\Gamma$ its chirality, and $A=D^2$. The classical angular spectrum
is $D=\pm n$, $n\ge1$, with multiplicity $2n$ per sign
\cite[(3.40), (3.54)]{CamporesiHiguchi1996}. Write $P_n$ for the rank-$4n$
projection of $A$ at $n^2$. With $P=-\ii\partial_v$ and $X=v$, put
\begin{equation}\label{eq:operators}
 \Dc=\Gamma P+2D,\qquad L_0=\Dc^2=P^2+4A,\qquad
 H_{\rm kin}=\varepsilon L_0,\quad
 \varepsilon=\frac{\hbar^2}{2m_{\rm phys}\ell_0^2},\quad
 \omega=\varepsilon/\hbar>0.
\end{equation}
The product connection and $\Gamma D=-D\Gamma$ remove mixed terms. On
$\dd v^2+g_{S^2}/4$, the scalar curvature is $8$; the Dirac square is
$(\nabla^{\mathcal C})^*\nabla^{\mathcal C}+2$. Thus the fixed curvature
contribution $2\varepsilon$ is already in \eqref{eq:operators}
\cite{TNDBackground,TNDParent}. No potential, wall or radial confinement
is inserted. Metric volume is $\dd v\dd\Omega/4$; multiplication by $2$
identifies product-measure spinors with that convention.

Use $\widehat\psi(k)=(2\pi)^{-1/2}\int\ee^{-\ii kv}\psi(v)\dd v$.
Choosing $\chi^-_{n,a}=\Gamma\chi^+_{n,a}$ makes each Fourier block
\begin{equation}\label{eq:blocks}
 M_n(k)=\begin{pmatrix}2n&k\\k&-2n\end{pmatrix},\quad
 M_n(k)^2=(k^2+4n^2)I,\quad
 \Pi_\pm(n,k)=\frac12\left(I\pm\frac{M_n(k)}{\sqrt{k^2+4n^2}}\right).
\end{equation}
There are $2n$ copies of each sign, all $n\ge1$, and all real $k$.
The square has multiplicity $4n$. The maximal domains are
\begin{align}
 \Dom L_0^p&=\left\{\psi\in\HH:
  \sum_n\int_\R(k^2+4n^2)^{2p}\norm{P_n\widehat\psi(k)}^2\dd k<\infty\right\},
 \label{eq:graphdomain}\\
 \Dom\Dc&=\left\{\psi\in\HH:
  \sum_n\int_\R(k^2+4n^2)\norm{P_n\widehat\psi(k)}^2\dd k<\infty\right\},
 \label{eq:Diracdomain}\\
 \Dom X^q&=\{\psi:|v|^q\psi\in L^2\},\qquad
 \Dom P^q=\{\psi:|k|^q\widehat\psi\in L^2\}.
 \label{eq:XPdomain}
\end{align}
Real Hermitian multiplication on these maximal domains gives the
self-adjoint realizations. Complete-shell cutoffs converge in graph norm;
for finitely many shells the graph norm is a finite sum of Sobolev norms,
so radial cutoff and mollification give a smooth core. In particular
$\Dom L_0=\Dom P^2\cap\Dom A$, because the nonnegative commuting
multipliers dominate their constituents and are bounded by their sum.
The radial spectrum is continuous. Generalized Fourier modes are not
normalizable states, and evenness reduces the square rather than $\Dc$.

\subsection{The source supplied by the background dictionary}
The full-line inverse of $-\partial_v^2+4n^2$ has kernel
$\ee^{-2n|v-w|}/(4n)$: its derivative has jump $-1$.
Its diagonal angular operator is $W=(4\sqrt A)^{-1}$, and
\begin{equation}\label{eq:Q}
 Q(\beta):=\Tr_\K(W\ee^{-\beta A})
 =\sum_{n\ge1}4n\frac{\ee^{-\beta n^2}}{4n}
 =\sum_{n\ge1}\ee^{-\beta n^2},\qquad\beta>0.
\end{equation}
All fixed $\beta$ derivatives converge locally uniformly. The prescribed
scan $\beta(v)=\pi\ee^{4v}$ and $b(v)=\ee^vQ(\beta(v))$ satisfy
\begin{equation}\label{eq:sourceconstruction}
 \Phi(v)=\frac18(b''(v)-b(v)).
\end{equation}
Differentiating the complete series yields \eqref{eq:Phi}. The Green diagonal
is position independent; the scan and imprint are supplied, not consequences
of translating a probe in the free geometry.

Let $\Theta(R)=1+2Q(\pi R^2)$. Jacobi inversion
$\Theta(R)=R^{-1}\Theta(1/R)$ follows from the principal-root specialization
of \cite[20.7.32]{DLMFThetaModular}. With $R=\ee^{2v}$ it gives
$b(v)-b(-v)=-\sinh v$. Applying $(\partial_v^2-1)/8$ proves smooth evenness
of the complete source. For $v\ge0$ every term of \eqref{eq:Phi} is positive,
since $\pi n^2\ee^{4v}>3/2$. Therefore $\Phi>0$ on the full line.
The zero mode in $\Theta$ is essential; finite reflected prefixes cannot
replace this smooth reflection argument.

\begin{lemma}[Complete source estimates]\label{lem:source}
Set $r=|v|$, $x=\pi\ee^{4r}$ and $\kappa=2-3/\pi>0$. For each fixed
$j\ge0$ there is a finite $B_j$ such that
\begin{equation}\label{eq:sourcebounds}
 \kappa\ee^r x^2\ee^{-x}\le\Phi(v),\qquad
 |\Phi^{(j)}(v)|\le B_j\ee^r x^{j+2}\ee^{-x}.
\end{equation}
All polynomial and compact-parameter exponential weights used in this
paper are integrable against these bounds, with complete source tails.
\end{lemma}
\begin{proof}
On $v\ge0$ define
\[
 \mathsf P_0(x)=2x^2-3x,\qquad
 \mathsf P_{j+1}(x)=(1-4x)\mathsf P_j(x)+4x\mathsf P_j'(x).
\]
Differentiation gives
$\Phi^{(j)}(v)=\ee^v\sum_{n\ge1}\ee^{-n^2x}\mathsf P_j(n^2x)$,
with $\deg\mathsf P_j=j+2$. If $a_j$ is its absolute coefficient sum, put
$C_d(N)=\sum_{n>N}n^{2d}\ee^{-\pi(n^2-1)}$. Then
\begin{equation}\label{eq:sourcetail}
 \ee^v\sum_{n>N}\ee^{-n^2x}|\mathsf P_j(n^2x)|
 \le a_jC_{j+2}(N)\ee^v x^{j+2}\ee^{-x},\qquad C_d(N)\longrightarrow0.
\end{equation}
This follows from $x\ge\pi>1$ and Gaussian summability. Taking $N=0$
gives $B_j=a_jC_{j+2}(0)$. The first source term gives the lower bound;
complete reflection transfers both bounds to $v<0$.
For finite $M,B\ge0$ and fixed $a,d\ge0$, substitution
$x=\pi\ee^{4|v|}$ shows that
\begin{equation}\label{eq:J}
 J_{a,d}(M,B)=\int_\R |v|^a\ee^{Mv^2+(B+1)|v|}x^d\ee^{-x}\dd v<\infty.
\end{equation}
The transformed integrand is a fixed power of $x$, powers of $\log(x/\pi)$,
and $\exp[M\log^2(x/\pi)/16-x]$. The term $-x$ dominates all the positive
terms. The same bound with $C_{j+2}(N)$ pays each integrated source tail.
\end{proof}

It follows, by dominated integration on complex parameter compacts, that
$H_\lambda(z)$ is jointly entire and
\begin{equation}\label{eq:Hmixed}
 \partial_\lambda^a\partial_z^bH_\lambda(z)
 =\frac12\int_\R v^{2a}(\ii v)^b\ee^{\lambda v^2}\Phi(v)\ee^{\ii zv}\dd v.
\end{equation}
This proves the Newman equation and all its fixed finite derivatives.
Complex $z$ is an analytic continuation, not a bounded measurement label.

\subsection{The normalized square-root state}
Put $F_\lambda(v)=\ee^{\lambda v^2}\Phi(v)$ and
\begin{equation}\label{eq:f}
 Z_\lambda=\int_\R F_\lambda(v)\dd v=2H_\lambda(0),\qquad
 f_\lambda=Z_\lambda^{-1/2}\ee^{\lambda v^2/2}\sqrt\Phi.
\end{equation}
On $|\lambda|\le M$, $0<Z_{-M}\le Z_\lambda\le Z_M<\infty$.
For instance an explicit lower bound is the integral of the lower estimate
in \eqref{eq:sourcebounds} over $0\le v\le1$, multiplied by
$2\ee^{-Mv^2}$. Define
\begin{equation}\label{eq:R}
 R=\frac{W\ee^{-\pi A}}{Q(\pi)}
   =\sum_{n\ge1}\frac{\ee^{-\pi n^2}}{4nQ(\pi)}P_n,\qquad
 \alpha_b:=\Tr(A^bR)=\frac{\sum_{n\ge1}n^{2b}\ee^{-\pi n^2}}{Q(\pi)}.
\end{equation}
Here $R$ is an angular density operator, $\Tr R=1$, and every fixed
$\alpha_b$ is finite. It is not a Gibbs state of the infinite cylinder.

\begin{proposition}[Admissibility of the complete preparation]\label{prop:state}
Each fixed Schwartz seminorm of $f_\lambda$ and each fixed mixed
parameter/spatial derivative is bounded on compact real $\lambda$ intervals.
For $\rho^I_\lambda=|f_\lambda\rangle\langle f_\lambda|\otimes R$ and every
fixed integer $p\ge0$, the quadratic-form sandwich is positive trace class,
with the exact payment
\begin{equation}\label{eq:fullmoment}
 \Tr(L_0^p\rho^I_\lambda L_0^p)
 =\sum_{j=0}^{2p}\binom{2p}{j}4^{2p-j}\alpha_{2p-j}
                                  \norm{P^jf_\lambda}_2^2<\infty.
\end{equation}
\end{proposition}
\begin{proof}
The matching exponents in \eqref{eq:sourcebounds} give
$|\Phi^{(\ell)}/\Phi|\le(B_\ell/\kappa)x^\ell$.
Set $a_\ell=\Phi^{(\ell)}/\Phi$ and $h=\sqrt\Phi$.
The identities $a_\ell'=a_{\ell+1}-a_\ell a_1$ yield the finite recurrence
\begin{equation}\label{eq:sqrtrecurrence}
 S_0=1,\qquad
 S_{j+1}=\tfrac12a_1S_j+
 \sum_{\ell=1}^j\frac{\partial S_j}{\partial a_\ell}
                          (a_{\ell+1}-a_\ell a_1),\qquad h^{(j)}=hS_j.
\end{equation}
Every monomial of $S_j$ has weighted degree $j$ if $a_\ell$ has weight
$\ell$. Hence
$|h^{(j)}|\le c_j\ee^{r/2}x^{j+1}\ee^{-x/2}$ for finite $c_j$.
The weakened exponent $-x/2$ still dominates every fixed polynomial,
Gaussian and logarithmic loss. Finite Leibniz differentiation therefore
bounds all spatial derivatives of $\ee^{\lambda v^2/2}h$ in Schwartz
seminorms. Parameter derivatives introduce only powers of $v^2/2$.
The integrals $Z_\lambda^{(a)}=\int v^{2a}F_\lambda\dd v$ and the positive
lower bound for $Z_\lambda$ control derivatives of $Z_\lambda^{-1/2}$.
Smooth positivity and complete reflection handle $v=0$.

For \eqref{eq:fullmoment}, expand $(k^2+4n^2)^{2p}$ in the full spectral
integral. Tonelli, Plancherel and \eqref{eq:R} give exactly the displayed
finite sum. This is the convergent positive sum of weighted ensemble
rank-one vectors; it does not assert that every vector in the density's
closed support belongs to $\Dom L_0^p$.
\end{proof}

All physical angular tails are the corresponding Gaussian series. For
fixed $j,N\ge0$ and $K>0$ the radial tail obeys
\begin{equation}\label{eq:radialtail}
 \int_{|k|>K}|k|^{2j}|\widehat f_\lambda(k)|^2\dd k
 \le K^{-2N}\norm{P^{j+N}f_\lambda}_2^2.
\end{equation}
Thus approximation pays the full spectra. Constants depend on the orders;
no order-uniform estimate follows. A positive Schwartz function does not in
general have a Schwartz square root: the ratio argument above is essential.
These are the preparation results of \cite[Section 3]{TNDDirect}, now in the
notation used for the count and collision calculus.
